\label{sec:exp}

\vspace{-0.15cm}
\paragraph{Datasets} We evaluate on four datasets ranging from 2D objects (MNIST~\cite{lecun1998gradient}), 3D objects (ModelNet40~\cite{wu20153d} rigid object, SHREC15~\cite{3dor.20151064} non-rigid object) to real 3D scenes (ScanNet~\cite{dai2017scannet}). Object classification is evaluated by accuracy.
%Part segmentation is evaluated by average intersection over union on all categories and parts.
Semantic scene labeling is evaluated by average voxel classification accuracy following \cite{dai2017scannet}. We list below the experiment setting for each dataset:
\begin{itemize}
    \item MNIST: Images of handwritten digits with 60k training and 10k testing samples.
    \item ModelNet40: CAD models of 40 categories (mostly man-made). We use the official split with 9,843 shapes for training and 2,468 for testing.
    \item SHREC15: 1200 shapes from 50 categories. Each category contains 24 shapes which are mostly organic ones with various poses such as horses, cats, etc. We use five fold cross validation to acquire classification accuracy on this dataset.
    %\item ShapeNetPart: This is a dataset containing 16,881 shapes from 16 classes, annotated with 50 parts in total. We use the official train test split following~\cite{shapenet2015}.
    \item ScanNet: 1513 scanned and reconstructed indoor scenes. We follow the experiment setting in~\cite{dai2017scannet} and use 1201 scenes for training, 312 scenes for test.
\end{itemize}


%\emph{Note: since it's not an application paper, include more self comparisons that show insights.}
\vspace{-0.3cm}
\subsection{Point Set Classification in Euclidean Metric Space}

%To classify an input point set, we first summarize the set into a global feature using our end-to-end trained network. Then a linear classifier with softmax loss is used for classification.
%In set classification task, given an input point set we want to classify it into one of the predefined categories. The set is firstly summarized into a global set feature and then a linear classifier with softmax loss is used for classification. The hierarchical feature learning network and the classifier is trained end-to-end.
We evaluate our network on classifying point clouds sampled from both 2D (MNIST) and 3D (ModleNet40) Euclidean spaces. MNIST images are converted to 2D point clouds of digit pixel locations. 3D point clouds are sampled from mesh surfaces from ModelNet40 shapes. In default we use 512 points for MNIST and 1024 points for ModelNet40. In last row (ours normal) in Table~\ref{tab:modelnet}, we use face normals as additional point features, where we also use more points ($N=5000$) to further boost performance. All point sets are normalized to be zero mean and within a unit ball. We use a three-level hierarchical network with three fully connected layers~\footnote{See supplementary for more details on network architecture and experiment preparation.}


%As a result, the input is $(x,y,z)$ coordinates and normals $(n_x, n_y, n_z)$.
%This turns out to be helpful in capturing fine-grained geometry structure of the objects. Also we increase $N$ to 5000 in this experiment.
%We also use more points in this experiment, resulting in an input size of $5000 \times 6$.


%\paragraph{Data preparation.} For MNIST images, we firstly normalize all pixel intensities to range $[0,1]$ and then select all pixels with intensities larger than 0.5 as valid digit pixels. Then we convert digit pixels in an image into a 2D point cloud with coordinates within $[-1,1]$, where the image center is the origin point. Augmented points are created to add the point set up to a fixed cardinality (512 in our case). We jitter the initial point cloud (with random translation of Gaussian distribution $\mathcal{N}(0,0.01)$ and clipped to 0.03) to generate the augmented points. For ModelNet40, we uniformly sample $N$ points from CAD models surfaces based on face area.
%To prepare the dataset, we first oversample on the mesh and then conduct farthest point sampling for a set of $N$ points. These $N$ points will be used for training and testing.
%By default, $N=1024$.
%For MNIST, the input is $(x,y)$ coordinates with a size $512\times2$. For ModelNet40, the input is $(x,y,z)$ coordinate with a size $1024 \times 3$.


%\paragraph{Network architecture and Training.} Following the protocols above, we use a network architecture with three set abstraction layers and two fully connected layers followed by a linear classifier with $K$ output scores ($K=10$ for MNIST, $K=40$ for ModelNet40):


%For all experiments, we use Adam~\cite{adam} optimizer with learning rate 0.001 for training. For data augmentation, we randomly scale object, perturb the object location as well as point sample locations. We also follow~\cite{qi2016volumetric} to randomly rotate objects for ModelNet40 data augmentation.
\vspace{-0.3cm}
\paragraph{Results.} In Table~\ref{tab:mnist} and Table~\ref{tab:modelnet}, we compare our method with a representative set of previous state of the arts. Note that PointNet (vanilla) in Table~\ref{tab:modelnet} is the the version in~\cite{qi2016pointnet} that does not use transformation networks, which is equivalent to our hierarchical net with only one level.

Firstly, our hierarchical learning architecture achieves significantly better performance than the non-hierarchical PointNet~\cite{qi2016pointnet}. In MNIST, we see a relative 60.8\% and 34.6\% error rate reduction from PointNet (vanilla) and PointNet to our method. In ModelNet40 classification, we also see that using same input data size (1024 points) and features (coordinates only), ours is remarkably stronger than PointNet.
Secondly, we observe that point set based method can even achieve better or similar performance as mature image CNNs. In MNIST, our method (based on 2D point set) is achieving an accuracy close to the Network in Network CNN. In ModelNet40, ours with normal information significantly outperforms previous state-of-the-art method MVCNN~\cite{su15mvcnn}. %\cite{qi2016volumetric} argues 3D based deep learning method on 3D data recognition was falling behind multi-view and CNN based method due to two issues: in-mature network architecture design and loss in data resolution (high-res image v.s. low-res voxel grid). Now given the hierarchical network design of PointNet++, as well as its ability of directly operating on high resolution point cloud data, we can achieve superior 3D feature learning and gain better performance than multi-view based methods.

\begin{table}[t!]
    \begin{minipage}{.45\textwidth}
        \centering
        \small
        \begin{tabular}{lc}
        \toprule
        Method     & Error rate (\%) \\
        \midrule
        Multi-layer perceptron~\cite{simard2003best} & 1.60     \\
        LeNet5~\cite{lecun1998gradient}     & 0.80      \\
        Network in Network~\cite{lin2013network}      & \textbf{0.47} \\
        %Kd-network~\cite{klokov2017escape}  & pc & 99.10 \\
        PointNet (vanilla)~\cite{qi2016pointnet}    & 1.30  \\
        PointNet~\cite{qi2016pointnet}    & 0.78  \\ \midrule
        Ours    & 0.51 \\
        \bottomrule
        \end{tabular}
        \caption{MNIST digit classification.}
        \label{tab:mnist}
    \end{minipage}
    \begin{minipage}{.55\textwidth}
        \centering
        \small
        \begin{tabular}[width=\linewidth]{lcccc}
        \toprule
        Method  & Input & Accuracy (\%) \\
        \midrule
        %VoxNet~\cite{maturana2015voxnet} &vox & 85.9 \\
        Subvolume~\cite{qi2016volumetric} &vox & 89.2 \\
        %O-CNN~\cite{wang2017cnn} &vox & 90.6 \\
        MVCNN~\cite{su15mvcnn} &img & 90.1 \\
        %Kd-network~\cite{klokov2017escape} &pc & 90.6 (91.8) \\
        PointNet (vanilla)~\cite{qi2016pointnet} &pc & 87.2 \\
        PointNet~\cite{qi2016pointnet} &pc & 89.2 \\ \midrule
        Ours  &pc & 90.7 \\ 
        Ours (with normal)   &pc&  \textbf{91.9} \\
        \bottomrule
        \end{tabular}
        \caption{ModelNet40 shape classification.}
        \label{tab:modelnet}
    \end{minipage}
  \end{table}

\vspace{-0.3cm}
\paragraph{Robustness to Sampling Density Variation.}
Sensor data directly captured from real world usually suffers from severe irregular sampling issues (Fig.~\ref{fig:rawscans}).
%Therefore, being robust to different sampling density is an important requirements for a point cloud processing network.
Our approach selects point neighborhood of multiple scales and learns to balance the descriptiveness and robustness by properly weighting them. %These make it very robust to sampling density variation.

\begin{figure}[t!]
\vspace{-0.5cm}
\centering
\begin{subfigure}
\centering
\includegraphics[width=0.45\linewidth,height=3cm]{./fig/chairs}
\end{subfigure}
\begin{subfigure}
\centering
\includegraphics[width=0.5\linewidth]{./fig/robustness3}
\end{subfigure}
\caption{Left: Point cloud with random point dropout. Right: Curve showing advantage of our density adaptive strategy in dealing with non-uniform density. DP means random input dropout during training; otherwise training is on uniformly dense points. See~Sec.\ref{sec:pointnet2++} for details.}
\label{fig:rb_classification}
\vspace{-0.5cm}
\end{figure}


%We verify this point in object classification task.
We randomly drop points (see Fig.~\ref{fig:rb_classification} left) during test time to validate our network's robustness to non-uniform and sparse data.
%, and an evaluation is shown on the right about how the performance of different methods would change.
%We evaluate our approach in four settings (SSG, SSG+DP, MSG+DP, MSG-XL+DP) and compare with two baseline method (PointNet vanilla, PointNet vanilla with DP).
%Here DP means the network is trained while randomly dropping out points in training sets, otherwise the network will be trained on relatively uniform and dense point sets.
%SSG means the network only groups points on a single scale in a set abstraction layer. MSG refers to multi-scale grouping and MSG-XL refers to cross-level multi-scale grouping.
In Fig.~\ref{fig:rb_classification} right, we see MSG+DP (multi-scale grouping with random input dropout during training) and MRG+DP (multi-resolution grouping with random input dropout during training) are very robust to sampling density variation. MSG+DP performance drops by less than $1\%$ from $1024$ to $256$ test points. Moreover, it achieves the best performance on almost all sampling densities compared with alternatives.
PointNet vanilla \cite{qi2016pointnet} is fairly robust under density variation due to its focus on global abstraction rather than fine details. However loss of details also makes it less powerful compared to our approach.
SSG (ablated PointNet++ with single scale grouping in each level) fails to generalize to sparse sampling density while SSG+DP amends the problem by randomly dropping out points in training time.
%MSG+DP and MSG-XL+DP go one step further, enabling the network to look at neighborhood of different scales during set abstraction. The network could then learn to balance different scales according to input sampling density and achieve both robustness and high descriptive power as a result.

\vspace{-0.3cm}
\subsection{Point Set Segmentation for Semantic Scene Labeling}
%\rqi{Message to convey: Our point set based method is comparable to state-of-the-art graph CNN method. Our hierarchical network is also better than original PointNet in segmentation task and semantic segmentation tasks, therefore empirically proving its effectiveness in local feature learning.}

% \begin{figure}[t!]
% \centering
% \begin{subfigure}
% \centering
% \includegraphics[width=0.55\linewidth]{./fig/sceneseg_larger_wmrg.pdf}
% \caption{Scannet labeling accuracy. The network is trained on ScanNet training scenes. Blue bar represents labeling accuracy on test scenes. Yellow bar represents accuracy on virtual scans synthesized from test scenes.}
% \end{subfigure}
% \begin{subfigure}
% \centering
% \includegraphics[width=0.4\linewidth]{./fig/sceneseg_vis.pdf}
% \end{subfigure}
% \caption{Scannet labeling results. \cite{qi2016pointnet} usually captures the overall layout of the room correctly but fails to discover the furniture. Our approach, in contrast, is much better at segmenting objects besides the room layout.}
% \label{fig:xxx}
% \end{figure}

%\todo{add network architecture parameter..}

\begin{wrapfigure}{R}{7cm}
  \vspace{-10pt}
  \begin{center}
    \includegraphics[width=7cm]{./fig/sceneseg_larger_wmrg.pdf}
  \end{center}  \vspace{-15pt}
  \caption{Scannet labeling accuracy. %The network is trained on ScanNet training scenes. Blue bar represents labeling accuracy on test scenes. Yellow bar represents accuracy on virtual scans synthesized from test scenes.% The virtual scans are non-uniformly sampled. Our approach outperforms baselines on both cases. Multi-scale grouping with dropout training achieves the best performance and best robustness regarding sampling density variation.
  }
  \label{fig:sceneseg}
\end{wrapfigure}

To validate that our approach is suitable for large scale point cloud analysis, we also evaluate on semantic scene labeling task. The goal is to predict semantic object label for points in indoor scans. \cite{dai2017scannet} provides a baseline using fully convolutional neural network on voxelized scans. They purely rely on scanning geometry instead of RGB information and report the accuracy on a per-voxel basis. To make a fair comparison, we remove RGB information in all our experiments and convert point cloud label prediction into voxel labeling following \cite{dai2017scannet}.
%We evaluate our framework in three settings, a basic version as is described in \ref{sec:pointnet2} (Ours1), adding drop out training (Ours2), network with adaptive scale selection described in \ref{sec:pointnet2++} (Ours3).
We also compare with \cite{qi2016pointnet}. %, a learning framework based on point cloud representation.
The accuracy is reported on a per-voxel basis in Fig.~\ref{fig:sceneseg} (blue bar).

Our approach outperforms all the baseline methods by a large margin. In comparison with \cite{dai2017scannet}, which learns on voxelized scans, we directly learn on point clouds to avoid additional quantization error, and conduct data dependent sampling to allow more effective learning. Compared with \cite{qi2016pointnet}, our approach introduces hierarchical feature learning and captures geometry features at different scales. This is very important for understanding scenes at multiple levels and labeling objects with various sizes.
%While evaluating on ScanNet scenes, which has been post-processed and re-sampled to make the sampling density uniform, SSG and MSG-DP do not have much performance gap. But still MSG-DP achieves the best performance.
We visualize example scene labeling results in Fig.~\ref{fig:sceneseg_vis}.
%Point cloud in ScanNet has been post-processed and re-sampled to make the sampling density uniform. Therefore the three settings of our framework do not differ much since the main difference in their design lies in coping with the irregular sampling issue.

\begin{wrapfigure}{R}{6.5cm}
  \vspace{-10pt}
  \begin{center}
    \includegraphics[width=6.5cm]{./fig/sceneseg_vis.pdf}
  \end{center}  \vspace{-15pt}
  \caption{Scannet labeling results. \cite{qi2016pointnet} captures the overall layout of the room correctly but fails to discover the furniture. Our approach, in contrast, is much better at segmenting objects besides the room layout.}
  \label{fig:sceneseg_vis}
  \vspace{-0.3cm}
\end{wrapfigure}

%We visualize some example scene labeling results in Figure~\ref{fig:sceneseg_vis}.
%\cite{qi2016pointnet} usually captures the overall layout of the room correctly but fails to discover the furniture. Our approach, in contrast, is much better at segmenting objects besides the room layout. % Focusing on evaluating the descriptive power of our network, we do not add additional post processing to enforce point labeling consistency, which could be achieved through a point cloud CRF as is done in \cite{Yi16}.

%\begin{table}[h!]
%\centering
%\begin{tabular}{@{}lccc}
%\toprule
%               & 3DCNN\cite{dai2017scannet} & PointNet\cite{qi2016pointnet} & Ours\\ \midrule
%Accuracy & 0.730 & 0.739 & \textbf{0.845}\\\bottomrule
%\end{tabular}
%\caption{Semantic voxel labeling accuracy on ScanNet test scenes.}
%\label{tab:sceneseg}
%\end{table}

\vspace{-0.3cm}
\paragraph{Robustness to Sampling Density Variation}
%We verify the robustness of our approach on semantic scene labeling task.
%Point cloud in ScanNet has been post-processed and re-sampled to make the sampling density uniform.
To test how our trained model performs on scans with non-uniform sampling density, we synthesize virtual scans of Scannet scenes similar to that in Fig.~\ref{fig:rawscans} and evaluate our network on this data. We refer readers to supplementary material for how we generate the virtual scans.
We evaluate our framework in three settings (SSG, MSG+DP, MRG+DP) and compare with a baseline approach \cite{qi2016pointnet}.
%We evaluate our framework in three settings, a basic version as is described in \ref{sec:pointnet2} (Ours1), adding drop out training (Ours2), network with adaptive scale selection described in \ref{sec:pointnet2++} (Ours3).

Performance comparison is shown in Fig.~\ref{fig:sceneseg} (yellow bar). We see that SSG performance greatly falls due to the sampling density shift from uniform point cloud to virtually scanned scenes. MRG network, on the other hand, is more robust to the sampling density shift since it is able to automatically switch to features depicting coarser granularity when the sampling is sparse. Even though there is a domain gap between training data (uniform points with random dropout) and scanned data with non-uniform density, our MSG network is only slightly affected and achieves the best accuracy among methods in comparison. These prove the effectiveness of our density adaptive layer design.
%Again, \cite{qi2016pointnet} does not perform very well in scenes due to its weakness in capturing multi-scale geometry features.
%SSG is trained on point cloud with monotonous sampling density and it does not consider how to select good neighborhood in non-uniformly sampled points. Therefore SSG does not generalize very well across different sampling densities.
%By randomly dropping points in the training point clouds as is done in MSG+DP, the network will have a chance to look at data with different sampling density, therefore better generalize to the virtual scans. Moreover, by adding adaptive scale selection mechanism, the network could learn to weight neighborhoods with different scales properly in different sampling density, thus achieving the best performance and the least performance drop.

%\begin{table}[h!]
%\centering
%\begin{tabular}{@{}cccc}
%\toprule
%               PointNet\cite{qi2016pointnet} & Ours1 & Ours2 & Ours3\\ \midrule
%0.545 & 0.719 & 0.739 & \textbf{0.751} \\ \bottomrule
%\end{tabular}
%\caption{Semantic voxel labeling accuracy on virtual scans generated from ScanNet test scenes. Ours1 corresponds to the basic version of our network without dropout training and adaptive scale selection. Ours2 is trained via randomly dropped out points in the training point clouds. Ours3, in addition, contains adaptive scale selection.}
%\label{tab:rb_sceneseg}
%\end{table}

%\begin{figure}[h]
%\centering
%\includegraphics[width=0.5\linewidth]{./fig/rb_sceneseg.png} 
%\caption{Using ScanNet as an example, showing how accuracy drops when sampling density decreases. Include PointNet1 and PointNet2, PointNet2 with drop out and PointNet with adaptive scale selection.}
%\label{fig:rb_sceneseg}
%\end{figure}

%We observe similar behavior on semantic scene labeling task, as is shown in Figure~\ref{fig:rb_sceneseg}. To emphasize the influence of object categories other than ``wall'', ``floor'', we use average per category point labeling accuracy instead of overall point labeling accuracy in Figure~\ref{fig:rb_sceneseg}.
\vspace{-0.3cm}
\subsection{Point Set Classification in Non-Euclidean Metric Space}

In this section, we show generalizability of our approach to non-Euclidean space. In non-rigid shape classification (Fig.~\ref{fig:intrinsic_shape_diff}), a good classifier should be able to classify $(a)$ and $(c)$ in Fig.~\ref{fig:intrinsic_shape_diff} correctly as the same category even given their difference in pose, which requires knowledge of intrinsic structure.
%This requires the ability to recognize intrinsic structure of an object, which is invariant to isometric deformation.
Shapes in SHREC15 are 2D surfaces embedded in 3D space. Geodesic distances along the surfaces naturally induce a metric space. We show through experiments that adopting PointNet++ in this metric space is an effective way to capture intrinsic structure of the underlying point set.

\begin{wrapfigure}{R}{4.5cm}
  \vspace{-15pt}
  \begin{center}
    \includegraphics[width=4.5cm]{./fig/intrinsic_shape_diff.png}
  \end{center}  \vspace{-15pt}
  \caption{An example of non-rigid shape classification.}
  \label{fig:intrinsic_shape_diff}
\end{wrapfigure}

For each shape in \cite{3dor.20151064}, we firstly construct the metric space induced by pairwise geodesic distances. We follow \cite{rustamov2009interior} to obtain an embedding metric that mimics geodesic distance. Next we extract intrinsic point features in this metric space including WKS \cite{aubry2011wave}, HKS \cite{sun2009concise} and multi-scale Gaussian curvature \cite{meyer2002discrete}. We use these features as input and then sample and group points according to the underlying metric space. In this way, our network learns to capture multi-scale intrinsic structure that is not influenced by the specific pose of a shape. Alternative design choices include using $XYZ$ coordinates as points feature or use Euclidean space $\mathbb{R}^3$ as the underlying metric space. We show below these are not optimal choices.

% Our approach sample and group the underlying point set based on the metric space it lies in. The metric space does not have to be an Euclidean space. In this section, we show the generalizability of our approach in non-Euclidean metric space. We consider the task of non-rigid shape classification and an example is shown in Figure~\ref{fig:intrinsic_shape_diff}. A good classifier should be able to classify $(a)$ and $(c)$ correctly as being the same category given their difference in pose. This requires the ability of considering the intrinsic structure of the underlying objects, which is invariant to the isometric deformation of a shape. The shapes we are considering could be regarded as 2D surfaces embedded in 3D space and geodesic distances along the surfaces would naturally induce a metric space. We show through experiments that adopting our PointNet++ in this metric space is an effective way to capture the intrinsic structure of the underlying point set.
% For each shape in \cite{3dor.20151064}, we first construct the metric space induced by pairwise geodesic distances on shape surfaces. We follow the construction in \cite{rustamov2009interior} to obtain a embedding metric that mimics geodesic distance. Then we extract intrinsic point features in this metric space including WKS \cite{aubry2011wave}, HKS \cite{sun2009concise} and multiscale
% Gaussian curvature \cite{meyer2002discrete}. We use these features as input to our PointNet++, and we sample and group points according to the underlying metric space. This way, our network could  learn to capture multi-scale intrinsic structure but not influenced by the specific pose of a shape. Some alternative design choices include using $xyz$ coordiniates of points as the network input feature, using Euclidean space $\mathbb{R}^3$ as the underlying metric space. We show below these are not optimal choices for the non-rigid shape classification task.

\vspace{-0.3cm}
\paragraph{Results.}  
We compare our methods with previous state-of-the-art method \cite{luciano2017deep} in Table~\ref{tab:nonrigid_classification}. \cite{luciano2017deep} extracts geodesic moments as shape features and use a stacked sparse autoencoder to digest these features to predict shape category.
Our approach using non-Euclidean metric space and intrinsic features achieves the best performance in all settings and outperforms \cite{luciano2017deep} by a large margin.

Comparing the first and second setting of our approach, we see intrinsic features are very important for non-rigid shape classification. %This is because different categories mainly differ with respect to their intrinsic structure.
$XYZ$ feature fails to reveal intrinsic structures and is greatly influenced by pose variation.
Comparing the second and third setting of our approach, we see using geodesic neighborhood is beneficial compared with Euclidean neighborhood. Euclidean neighborhood might include points far away on surfaces and this neighborhood could change dramatically when shape affords non-rigid deformation. This introduces difficulty for effective weight sharing since the local structure could become combinatorially complicated. Geodesic neighborhood on surfaces, on the other hand, gets rid of this issue and improves the learning effectiveness.
%Thanks to the flexibility of PointNet++, we are able to process points in non-Euclidean metric space, which is very important for tasks like non-rigid shape classification shown here.

\begin{table}[h!]
\centering
\small
\begin{tabular}{lccc}
\toprule
& Metric space & Input feature & Accuracy (\%) \\ \midrule
DeepGM \cite{luciano2017deep} & - & Intrinsic features & 93.03 \\ \midrule
\multirow{3}{*}{Ours} & Euclidean & XYZ & 60.18\\
& Euclidean & Intrinsic features & 94.49\\
& Non-Euclidean & Intrinsic features & \textbf{96.09}\\\bottomrule
\end{tabular}
\caption{SHREC15 Non-rigid shape classification.}
\label{tab:nonrigid_classification}
\vspace{-0.7cm}
\end{table}

\subsection{Feature Visualization.}

\begin{wrapfigure}{R}{6cm}
  \vspace{-15pt}
  \begin{center}
    \includegraphics[width=6cm]{./fig/features.pdf}
  \end{center} 
  \caption{3D point cloud patterns learned from the first layer kernels. The model is trained for ModelNet40 shape classification (20 out of the 128 kernels are randomly selected). Color indicates point depth (red is near, blue is far). } \vspace{-25pt}
  \label{fig:visu}
  \vspace{-0.6cm}
\end{wrapfigure}

In Fig.~\ref{fig:visu} we visualize what has been learned by the first level kernels of our hierarchical network. We created a voxel grid in space and aggregate local point sets that activate certain neurons the most in grid cells (highest 100 examples are used). Grid cells with high votes are kept and converted back to 3D point clouds, which represents the pattern that neuron recognizes. Since the model is trained on ModelNet40 which is mostly consisted of furniture, we see structures of planes, double planes, lines, corners etc. in the visualization.

% \begin{figure}
% \centering
% \includegraphics[width=0.8\linewidth]{./fig/features}
% \caption{3D point cloud patterns learned from the first layer kernels (32 out of the 128 kernels are randomly selected). Color indicates point depth (red is near, blue is far). The model is trained for ModelNet40 shape classification.}
% \label{fig:visu}
% \end{figure}



% \begin{minipage}{\linewidth}
%       \centering
%       \begin{minipage}{0.45\linewidth}
%           \begin{figure}[H]
%               \includegraphics[width=\linewidth]{./fig/sceneseg_larger_wmrg}
%               \caption{This is the first figure}
%           \end{figure}
%       \end{minipage}
%       \hspace{0.05\linewidth}
%       \begin{minipage}{0.45\linewidth}
%           \begin{figure}[H]
%               \includegraphics[width=\linewidth]{./fig/sceneseg_vis}
%               \caption{This is the second figure}
%           \end{figure}
%       \end{minipage}
%   \end{minipage}
  